\documentclass[11pt]{article}
\usepackage[utf8]{inputenc}
\usepackage{amsmath}
\usepackage{booktabs}

% simple isotope helper (avoids the isotope package dependency)
\newcommand{\iso}[2]{\textsuperscript{#1}#2}

\begin{document}

\begin{center}
  {\large\textbf{PET-based proton range verification:}}\\[2pt]
  {\large\textbf{method, limitations, and the case for an early, short window}}\\[4pt]
  {\normalsize a reading of Parodi et al.\ (2007), Knopf et al.\ (2009),}\\[1pt]
  {\normalsize Espa\~na et al.\ (2011) and Grogg et al.\ (2013)}
\end{center}
\medskip

\noindent
This note summarises four Massachusetts General Hospital (MGH) studies of
PET-based range verification of proton therapy and draws out their common
prescription. The method compares the $\beta^+$-activity distribution produced
by the beam against a prediction, and reads the proton range from the distal
fall-off of that activity. The four studies span its development and its two
acquisition regimes --- \emph{offline} (a commercial PET/CT tens of minutes
away, tens of minutes after treatment) and \emph{in-room} (a mobile scanner at
the couch, acquiring within minutes):
\begin{itemize}
\item Parodi et al.~\cite{parodi2007a} (Sec.~1) --- the controlled offline
  phantom feasibility study: millimetre range accuracy under known geometry and
  composition.
\item Knopf et al.~\cite{knopf2009} (Secs.~2--7) --- the systematic offline
  patient-population study: in vivo, that accuracy is realised only in
  well-co-registered bony structures of head-and-neck patients and fails at the
  majority of positions and sites (washout, motion, HU mapping, geometry).
\item Espa\~na et al.~\cite{espana2011} (Sec.~8) --- the cross-section
  reliability study on an in-room CsI scanner: the achievable \emph{model}
  accuracy depends strongly on which isotope carries the signal, and hence on
  \emph{when} it is acquired.
\item Grogg et al.~\cite{grogg2013} (Sec.~9) --- a fast, MC-free in-room method
  comparing the distal \emph{endpoints} of dose and PET directly, which requires
  \iso{15}{O} dominance for a clean, single-threshold endpoint.
\end{itemize}
Taken together (Sec.~10) they argue for an \textbf{early-start, short-window}
acquisition on the \iso{15}{O} channel, for three independent reasons ---
statistical, cross-section-systematic, and endpoint-definition --- which is the
regime the present $\sigma_R$-vs-dose study operates in.

\subsection*{1. Phantom feasibility (Parodi et al.\ 2007)}

Parodi et al.~\cite{parodi2007a} irradiated plastic phantoms with SOBP proton
fields (15--16~cm water-equivalent range, 10~cm modulation, $8\times8$~cm$^2$
aperture) from the MGH 230~MeV cyclotron and imaged them 14--20~min later on an
LSO-based Biograph Sensation~16 in list mode, comparing with FLUKA/Geant4
predictions convolved with a 7~mm (transaxial) / 7.5~mm (axial) Gaussian. Four
configurations were used: a homogeneous PMMA cylinder (8~Gy single field, and
8~Gy two orthogonal fields), an inhomogeneous PMMA/bone/lung stack (2~Gy), and
PMMA with two titanium rods mimicking spinal implants (8~Gy). Three results
carry over to the present work.

\emph{Range accuracy.} The 50\% distal fall-off of the measured activity agreed
with the simulation \textbf{within 1~mm in every configuration} --- homogeneous
PMMA, the inhomogeneous stack, and even through the titanium implants --- and
held to within 1~mm when only the last 20~min was reconstructed (the
$\sim$2~Gy-equivalent low-statistics case, activity reduced by $2.29\times$).
Lateral field position was likewise verified to $<1$~mm. For the overlapping
two-field case the analysis had to be restricted to the 30--10\% fall-off to
exclude the second beam, and even there stayed within 1~mm.

\emph{Isotope content and yield.} A maximum \iso{11}{C} production of
$\sim$870~Bq\,ml$^{-1}$\,Gy$^{-1}$ was measured in PMMA after $\sim$6~min
irradiation. At the start of imaging the offline signal is $94$--$97\%$
\iso{11}{C}, $3$--$5\%$ \iso{13}{N}, with only trace \iso{15}{O} in the first
minutes --- the long-lived-isotope regime that offline acquisition selects.
Rescaled to tissue by carbon content, this corresponds to $\sim$0.3
(muscle)--1.4 (bone/adipose)~kBq\,ml$^{-1}$ after a 2~Gy patient treatment.

\emph{LSO intrinsic background.} The \iso{176}{Lu} content of LSO (2.6\% of
natural Lu) produces a \textbf{stationary $\sim$930--1000~cps random background}
independent of the signal. Because random coincidences scale as the square of the
singles rate, this intrinsic noise dominates at the low count rates of therapy
monitoring: the noise-to-signal ratio exceeded $0.2$ for true rates below
$\sim$4.8--5.7~kcps (about 1~kBq\,ml$^{-1}$, $\sim$1~MBq total). This is a
concrete penalty of an intrinsically radioactive scintillator in this
low-signal application --- the axis on which BGO and CsI, which carry no such
intrinsic activity, differ from LSO/LYSO. A separate finding, that PET range
verification is insensitive to CT artifacts from the titanium rods yet more
sensitive to the fluence ``shadowing'' behind them than a commercial pencil-beam
plan, is not pursued here.

\subsection*{2. Method and range metrics}

Range verification compares depth profiles of the measured and simulated
activity in the distal fall-off, at three definitions the authors label A, B, C:
\emph{method A} at the 20\% level of the last local maximum, \emph{method B} at
the 50\% level, and \emph{method C} (``shift method'') using the whole fall-off
region rather than a single point. The shift method is the more robust of the
three because it relies on the overall shape agreement rather than on a single
landmark, and it dominates the soft-tissue conclusions below.

\subsection*{3. Reproducibility (the enabling result)}

Two patients were imaged twice, one week apart, for the same treatment fraction.
The measured activity \emph{range} reproduced to $\sim$1~mm standard deviation
(mean absolute range deviation 1.14~mm and 0.125~mm for the two patients, shift
method), consistent with the 1~mm reproducibility seen earlier in phantoms.
Absolute activity values differed by up to $\pm30\%$ between the two scans, but
this is \emph{not} a timing effect: the MC, propagating the two different
irradiation durations and delays, predicts at most 3\% difference. The $\pm30\%$
therefore reflects statistical noise plus genuine day-to-day changes in tissue
composition and perfusion. \textbf{The 1~mm range reproducibility is what makes
the whole approach viable}; everything below is about the systematic offsets that
sit on top of it.

\subsection*{4. Biological washout}

Six head-and-neck patients (chosen for good immobilisation and rigid geometry, so
motion and co-registration are minimised) were analysed at positions sorted by
where the beam stops. The mean measured-vs-simulated range agreement is given in
Table~\ref{tab:washout}.

\begin{table}[h]
  \centering
  \caption{Mean agreement between measured and MC-simulated activity range,
    for the three stopping cases and the three range metrics
    (Knopf et al.~\cite{knopf2009}, Table~1). Case (i): beam stops in bone;
    case (ii): stops in soft tissue with the last activity maximum in bone;
    case (iii): stops in soft tissue with the last maximum also in soft tissue.}
  \label{tab:washout}
  \begin{tabular}{lccc}
    \toprule
    & \multicolumn{2}{c}{Pointwise (mm)} & Shift (mm) \\
    \cmidrule(lr){2-3}
    Case & 50\% (B) & 20\% (A) & (C) \\
    \midrule
    (i)   Bone                 & 2.5 & 1.2 & 2.4 \\
    (ii)  Bone / soft tissue   & 3.2 & 8.1 & 2.2 \\
    (iii) Soft tissue          & 6.8 & 3.9 & 4.3 \\
    \bottomrule
  \end{tabular}
\end{table}

\noindent
In bone (case i) all three metrics agree within $\sim$2.5~mm and verification
works, though individual positions can still deviate by up to 6.6~mm (50\%) /
3.6~mm (20\%). In soft tissue, pointwise verification breaks down because blood
perfusion degrades the measured distribution: the robust shift method reaches
4.3~mm in the worst case (iii) and 2.2~mm when the last maximum is anchored in
bone (case ii), but $27\%$ (case ii) / $33\%$ (case iii) of soft-tissue positions
are off by more than $\pm5$~mm. \textbf{Reliable mm-accurate verification in soft
tissue is therefore not feasible with the offline method.} The paper's often-quoted
``$\sim$4~mm washout discrepancy'' is precisely this case-(iii) soft-tissue number.
Two points matter for our reading: (a) this is a \emph{measured-vs-model bias}
driven by an imperfect washout model, not the intrinsic statistical $\sigma_R$;
and (b) the authors' proposed cure --- shorter delay via in-room PET to minimise
perfusion washout --- is the premise of the present scanner study.

\subsection*{5. Patient motion}

During the 30~min acquisition, breathing blurs the activity distribution. Lateral
conformity between measured and simulated distributions was up to 5~mm in
head-and-neck sites but degraded to as much as 30~mm in abdominopelvic sites,
dominant in the anterior--posterior direction. This excludes most thoracic and
abdominal tumours from mm-accurate verification with the 3D offline strategy,
independent of washout.

\subsection*{6. Uncertainty in the simulated activity (HU mapping)}

The MC activity depends on assigning elemental composition and washout parameters
to each voxel from its Hounsfield unit. Where tissues share an HU domain the
mapping is not correlative: bone marrow and soft tissue both fall in
$-20$ to $110$~HU, yet the Schneider conversion underestimates the marrow carbon
fraction by nearly 50\%, and since activation is dominated by the $(p,pn)$ channel
on \iso{12}{C} this underestimates marrow activity. Varying the HU conversion and
the washout parameters within plausible ranges moves the \emph{absolute} activity
by up to $\sim$35\% and the \emph{range} by $\sim$1.8--2.4~mm
(Table~\ref{tab:hu}), while the simulated \emph{dose} profiles are essentially
unaffected --- activity is far more sensitive to tissue parameters than dose is.

\begin{table}[h]
  \centering
  \caption{Activity-range differences induced by parameter changes in the MC
    (Knopf et al.~\cite{knopf2009}, Table~4): original vs adjusted HU conversion,
    and original vs adjusted washout parameters.}
  \label{tab:hu}
  \begin{tabular}{lccc}
    \toprule
    & 20\% (A) & 50\% (B) & Shift (C) \\
    \midrule
    $R_{\text{orig}}-R_{\text{adj.\ HU}}$ (mm)      & 1.80 & 2.41 & 1.81 \\
    $R_{\text{orig}}-R_{\text{adj.\ washout}}$ (mm) & 1.70 & 2.28 & 1.81 \\
    \bottomrule
  \end{tabular}
\end{table}

\noindent
The washout model itself is the space- and time-dependent weighting
$C_{\mathrm{bio}}(\mathbf{r},t) = M_s(\mathbf{r})\,
\exp[-\lambda_{\mathrm{bio}}(\mathbf{r})\,t]$ of the physical activity, with an
HU-based classification into low-, intermediate- and normally perfused tissue
(muscle: $M_s=0.55$, $T_{1/2,\mathrm{bio}}=3500$~s; brain: $M_s=0.35$,
$T_{1/2,\mathrm{bio}}=10\,000$~s), following Mizuno et al.~\cite{mizuno2003} and
neglecting any distinction between isotopes. These are the source constants used
in our own washout note.

\subsection*{7. Site-specific limitations}

Beyond the three general factors, offline verification is compromised for
particular sites by: \emph{beam arrangement} (opposed or small-angle fields blur
the distal edge of each field together --- feasible only for single or nearly
orthogonal beams); \emph{co-registration} (rigid CT fusion needs a rigid,
fully-imaged anatomy, which head/neck provides and eye/abdominopelvic do not);
and \emph{organ motion} (e.g.\ prostate patients voiding the bladder between
treatment and imaging changes the geometry irrecoverably).

\subsection*{8. Cross-section reliability: the case for acquiring early and
briefly (Espa\~na et al.\ 2011)}

The Knopf limitations above are dominated by patient physiology. Espa\~na et
al.~\cite{espana2011} isolate a different, purely nuclear source of
model error and, in doing so, give the strongest argument in this set for a
short, early acquisition window. They separate the cross-section input from the
transport engine: proton transport uses a Geant4-based code, but the
$\beta^+$-production \emph{cross sections} are supplied as external tabulated
data --- experimental values from EXFOR and theoretical values from ICRU
Report~63 --- convolved with the voxel proton fluence, so the cross-section
choice becomes the controlled variable rather than a fixed physics-list output.
Measurements used the CsI-based NeuroPET scanner (31~cm bore, 24~cm axial FOV)
and a three-material phantom (HDPE, gel-water, tissue-equivalent gel) that
isolates the three channels producing $\sim$95\% of the activity:
\iso{16}{O}$(p,pn)$\iso{15}{O}, \iso{12}{C}$(p,pn)$\iso{11}{C}, and
\iso{16}{O}$(p,3p3n)$\iso{11}{C}.

The decisive result is that the achievable measured-vs-model range agreement
depends strongly on the acquisition protocol, because the protocol selects which
isotope --- and hence which cross section --- carries the signal. In the
tissue-equivalent gel, with the best available cross sections combined, the
50\% distal-fall-off agreement was:

\begin{table}[h]
  \centering
  \caption{Measured-vs-Monte-Carlo range difference at the 50\% distal
    fall-off in tissue-equivalent gel, for the two acquisition protocols and
    both field types (Espa\~na et al.~\cite{espana2011}). The 5~min protocol
    reconstructs the first 5~min; the 30~min protocol skips 15~min and
    reconstructs the next 30~min.}
  \label{tab:espana}
  \begin{tabular}{lcc}
    \toprule
    Protocol & Monoenergetic (mm) & SOBP (mm) \\
    \midrule
    5~min (early)   & $0.2 \pm 0.5$ & $0.8 \pm 0.5$ \\
    30~min (delayed) & $2.1 \pm 0.5$ & $4.4 \pm 0.5$ \\
    \bottomrule
  \end{tabular}
\end{table}

\noindent
The mechanism is isotope weighting in time. Early, the signal is dominated by
\iso{15}{O} (half-life 122~s) from \iso{16}{O}$(p,pn)$\iso{15}{O}, whose cross
sections are comparatively well measured --- so the model bias is sub-millimetre.
By the delayed window \iso{15}{O} has decayed and the signal is carried by
\iso{11}{C}, for which the \iso{16}{O}$(p,3p3n)$\iso{11}{C} channel is poorly
known (experimental and ICRU values disagree by up to 50\% and produce a
spurious double fall-off), leaving a residual $\sim$4~mm even with the best data.
Secondary-neutron production channels contribute only 0.6\%. This is a
\emph{systematic} model bias, distinct from the statistical $\sigma_R$ our study
reports and from the physiological washout bias of Sec.~4: it argues for
acquiring \textbf{early} (while \iso{15}{O} is present) and \textbf{stopping
early} (before the \iso{11}{C}-dominated, cross-section-limited tail
accumulates). It also complements our own generation-2 finding by an independent
route: where our per-isotope study showed \iso{15}{O} is more \emph{statistically}
precise per count than \iso{11}{C}, Espa\~na et al.\ show it is also more
\emph{systematically} reliable --- both axes favour the short, early window.

\subsection*{9. A fast, MC-free endpoint method (Grogg et al.\ 2013)}

Grogg et al.~\cite{grogg2013} target the cost of the method rather than its
accuracy: a full Geant4 MC-PET takes $\sim$12~h on 10 dual-core CPUs per patient,
so they propose a fast screen ($<$30~min on one CPU) that compares the distal
\emph{endpoints} of the planned dose and the PET activity directly, with no
MC-PET and no filter functions. The rationale is the one our reconstruction
relies on: cross sections, elemental composition and washout all affect the
\emph{shape} of the activity profile, but its distal \emph{endpoint} is fixed by
the nuclear-reaction energy thresholds alone, so it maps to the end of range
independently of those uncertainties. The PET endpoint is defined as the
\textbf{x-intercept of the best linear fit to the distal fall-off} (fit window
from the last distal maximum extending 2.5~cm, sub-range chosen by best
residual-sum-of-squares over $\ge 8$ voxels); the dose endpoint is the midpoint
of its 20--50\%-of-maximum fall-off. They deliberately reject the single-point
``$X\%$ of the distal maximum'' definition as unreliable and ambiguous when the
profile carries a small spurious distal peak. A fixed intrinsic offset of
$1.4 \pm 0.1$~mm (PET endpoint shallower than dose, modulation-independent to
$<0.08$~mm) is measured in a gel phantom and subtracted.

Method and data are in-room: proton transport is the Geant4-based MGH system
with EXFOR/ICRU cross sections optimised as in Espa\~na et al.; the scanner is
the CsI-based mobile NeuroPET (patient results use simulated PET at 7~mm FWHM;
only a rabbit-thigh irradiation is real NeuroPET data), with an integrated
NeuroPET/CT (2.5~mm resolution) named as the platform for patient measurement.
Across six head-and-neck patients the mean dose$-$simPET endpoint difference was
$-1.6$ to $+0.5$~mm between patients, with a per-patient spread of $1.1$--$4.0$~mm
(Table~\ref{tab:grogg}); the rabbit gave $0.2 \pm 2.5$~mm. Realistic accuracy is
$\sim$3~mm with simPET-quality data and $\sim$4--5~mm with current real PET ---
short of millimetre, but enough to flag beams outside the MGH safety margin
($3.5\% + 1$~mm), with heterogeneous nasal-cavity paths driving the large-spread
cases.

\begin{table}[h]
  \centering
  \caption{Mean and standard deviation of the dose$-$simPET distal-endpoint
    difference for the six head-and-neck patients (Grogg et
    al.~\cite{grogg2013}, Table~II).}
  \label{tab:grogg}
  \begin{tabular}{lcccccc}
    \toprule
    Patient & A & B & C & D & E & F \\
    \midrule
    Mean (mm)    & $-1.6$ & $0.5$ & $-0.4$ & $0.4$ & $-0.9$ & $-1.7$ \\
    Std.\ dev.\ (mm) & $1.9$ & $1.6$ & $1.1$ & $2.7$ & $4.0$ & $2.5$ \\
    \bottomrule
  \end{tabular}
\end{table}

\noindent
The result that matters here is the requirement, not the accuracy. The endpoint
map is clean only when a \emph{single} nuclear threshold sets the fall-off, which
holds for \iso{15}{O} --- produced through the one channel
\iso{16}{O}$(p,pn)$\iso{15}{O} (threshold 16.79~MeV) --- but not for \iso{11}{C},
which comes from three channels of differing threshold that smear the endpoint.
Grogg et al.\ therefore run a 5~min in-room scan started 2~min after treatment,
when \iso{15}{O} is $\sim$80\% of the signal. This is a \emph{third}, independent
reason to prefer the early, \iso{15}{O}-dominated window, on top of the
statistical (per-count precision) and cross-section-systematic (Sec.~8) reasons.
Two further contacts with our work: they apply a reverse washout correction using
the same Mizuno carbon-beam model our washout note adopts, and flag it as
``questionable'' for protons --- the caveat our note also carries; and they are
developing a direct \iso{15}{O} kinetic model to remove the washout correction
entirely, which is the spatial-perfusion direction our washout note leaves open.

\subsection*{10. Conclusion and relevance here}

In phantoms, where geometry and composition are known and motion and
co-registration are absent, the offline method reaches 1~mm range accuracy even
at 2~Gy-equivalent statistics (Sec.~1). In vivo, despite the same 1~mm range
reproducibility, offline PET/CT verifies the proton range to 1--2~mm only in
low-perfused, well-co-registered bony structures of head-and-neck patients.
Washout limits \emph{where} within a field verification is possible ($\sim$4~mm
in soft tissue); motion limits \emph{which whole sites} qualify (up to 3~cm in
the abdomen); and the MC activity is itself limited to $\sim$35\% in absolute
value and $\sim$2~mm in range by HU-to-tissue mapping. The authors argue all of
these are ``conquerable'', principally by \textbf{in-room PET} --- reducing the
treatment-to-scan gap to $\sim$1~min (their movable NeuroPET) suppresses washout
and motion, and removes the repositioning that drives co-registration error.
Espa\~na et al.\ (Sec.~8) and Grogg et al.\ (Sec.~9) sharpen the timing
prescription from the opposite end: not only \emph{start} early to beat washout,
but keep the window \emph{short} so the signal stays on the \iso{15}{O} channel
rather than drifting onto the \iso{11}{C} channels. Three independent arguments
converge on that \iso{15}{O}, \textbf{early-start, short-window} acquisition:
(i) \emph{statistical} --- our own generation-2 study finds \iso{15}{O} more
precise per count; (ii) \emph{cross-section-systematic} --- \iso{15}{O}'s single
production channel is well measured, whereas the \iso{11}{C} channels are not
($\sim$4~mm residual model bias, Sec.~8); and (iii) \emph{endpoint-definition}
--- \iso{15}{O}'s single nuclear threshold gives one clean distal endpoint, while
\iso{11}{C}'s three thresholds smear it (Sec.~9). That prescription --- short
delay, short window, in-room acquisition on \iso{15}{O} --- is exactly the regime
the present $\sigma_R$-vs-dose study operates in. Two further threads connect
directly to our toolchain: the LSO intrinsic-background result (Sec.~1) is a
quoted argument on the scintillator-choice axis that separates BGO and CsI from
LSO/LYSO; and the Mizuno washout parameterisation used across these papers
(Secs.~6 and~9) is the one our downstream washout model adopts, complete with the
same caveat about its proton applicability.

\begin{thebibliography}{9}

\bibitem{parodi2007a}
K.~Parodi, H.~Paganetti, E.~Cascio, J.~B.~Flanz, A.~A.~Bonab, N.~M.~Alpert,
K.~Lohmann and T.~Bortfeld,
``PET/CT imaging for treatment verification after proton therapy: a study with
plastic phantoms and metallic implants,''
\textit{Med.\ Phys.} \textbf{34} (2007) 419--435.

\bibitem{knopf2009}
A.~Knopf, K.~Parodi, T.~Bortfeld, H.~A.~Shih and H.~Paganetti,
``Systematic analysis of biological and physical limitations of proton beam
range verification with offline PET/CT scans,''
\textit{Phys.\ Med.\ Biol.} \textbf{54} (2009) 4477--4495.

\bibitem{espana2011}
S.~Espa\~na, X.~Zhu, J.~Daartz, G.~El~Fakhri, T.~Bortfeld and H.~Paganetti,
``Reliability of proton-nuclear interaction cross-section data to predict
proton-induced PET images in proton therapy,''
\textit{Phys.\ Med.\ Biol.} \textbf{56} (2011) 2687--2698.

\bibitem{grogg2013}
K.~Grogg, X.~Zhu, C.~H.~Min, B.~Winey, T.~Bortfeld, H.~Paganetti, H.~A.~Shih and
G.~El~Fakhri,
``Feasibility of using distal endpoints for in-room PET range verification of
proton therapy,''
\textit{IEEE Trans.\ Nucl.\ Sci.} \textbf{60} (2013) 3290--3297.

\bibitem{mizuno2003}
H.~Mizuno, T.~Tomitani, M.~Kanazawa et al.,
``Washout measurements of radioisotope implanted by radioactive beams in the
rabbit,''
\textit{Phys.\ Med.\ Biol.} \textbf{48} (2003) 2269--2281.

\end{thebibliography}

\end{document}
